\documentclass[14pt]{extarticle}
% math symbols
\usepackage{sfg}


\usepackage{amssymb,amsmath}
\synctex=1
% for different compilers
\usepackage{ifpdf}
% geometry of page
\usepackage[margin=2cm]{geometry}

% if pdflatex, then
\ifpdf
\usepackage[russian]{babel}
\usepackage[utf8]{inputenc}
\usepackage[unicode]{hyperref}
\usepackage[pdftex]{graphicx}
\usepackage{cmlgc}
% if xelatex, then
\else
% math fonts
\usepackage{fouriernc}
% xelatex specific packages
\usepackage[xetex]{hyperref}
\usepackage{xltxtra}	% \XeLaTeX macro
\usepackage{xunicode}	% some extra unicode support
\defaultfontfeatures{Mapping=tex-text}
\usepackage{polyglossia}	% instead of babel in xelatex
\usepackage{indentfirst}	% 
\setdefaultlanguage{russian}
% fonts
\newfontfamily\cyrillicfont{SchoolBookC}
\newfontfamily\cyrillicfontsf{TextBookC}
\setmonofont{Consolas}
\fi

% several pictures in one figure
\usepackage{subfig}
% calc in TeX expressions
\usepackage{calc}
% nice pictures and plots
\usepackage{pgfplots,tikz,circuitikz}
% different libraries for pictures
\usetikzlibrary{%
  arrows,%
  calc,%
  patterns,%
  decorations.pathreplacing,%
  decorations.pathmorphing,%
  decorations.markings,%
  intersections,%
  decorations.text%
}

\usepackage{tkz-euclide}

\usepackage{enumitem}
\renewcommand{\theenumi}{(\asbuk{enumi})}
\renewcommand{\labelenumi}{\asbuk{enumi})}
\AddEnumerateCounter{\Asbuk}{\@Asbuk}{\CYRM}
\AddEnumerateCounter{\asbuk}{\@asbuk}{\cyrm}

\begin{document}

\section*{Задача}

\subsection*{Условие}
Пусть $A\in\alpha$ и прямая $AB$ перпендикулярна двум прямым $AC$ и $AD$ плоскости $\alpha$. Проведите прямую $AK$ в плоскости $\alpha$ и возьмите на ней любую точку $X\neq A$. Через точку X проведите отрезок, который заключен между прямыми $AC$ и $AD$ и точкой $X$ делится пополам. Соедините точку B с концами этого отрезка и токой $X$. Исходя из этого построения докажите, что $AB\perp\alpha$ 
\subsection*{Решение}
Так как $AB\perp AC$ и $AB\perp AD$, $AB\perp\alpha$. Из формулы медианы

\begin{equation}
|AX|^2=\frac{2|AC|^2|AD|^2-|CD|}{4}.
\end{equation}

\begin{equation}
|BX|^2=\frac{2|BC|^2|BD|^2-|CD|}{4}.
\end{equation}

А значит, что

\begin{equation}
|BX|^2-|AX|^2 = |AB|^2.
\end{equation}

Из обратной теоремы Пифагора получаем, что $AB\perp \alpha$.
\newpage
\begin{figure}[h]
	\centering
	\includegraphics[width=1\textwidth]{{/home/galqiwi/geoma/7.25}.png}
\end{figure}


\end{document}